\section{Scope, System Model, and Problem Formulation}

\subsection{Scope and assumptions}

The confirmatory experiment is a \emph{source-partitioned inference simulation}, not a physical 5G deployment. 5G-NIDD provides source-side metadata, denoted \texttt{sVid}, that is excluded from predictor training and used only to assign flows to eight monitor slots through a stable hash. Each flow has one owner source. The principal method transmits one owner message to the coordinator; it does not claim consensus among simultaneously observing agents. Accordingly, we use the terms ``source monitor'' and ``multi-source'' throughout.

The reported bytes are application-layer evidence payloads. They exclude lower-layer IP, UDP, QUIC, Message Queuing Telemetry Transport (MQTT), Constrained Application Protocol (CoAP), and 5G user-plane headers. The message can be carried by existing transports; no radio-standard modification is required \cite{rfc7252,mqtt5}. Network latency, retransmission, authentication, and scheduler interactions are outside the main benchmark and are examined only in exploratory stress tests.

The threat model assumes attacks alter flow statistics but do not initially compromise the source monitor. A separate experiment manipulates class, attribution, and anomaly fields or replays benign messages. Those attacks expose vulnerabilities; they are not presented as a solved Byzantine-resilience problem.

\subsection{Notation}

Let $\mathcal{Y}=\{1,\ldots,K\}$ be the known label set, and let $h\notin\mathcal{Y}$ be the family withheld for open-set testing. After leakage-controlled preprocessing, a flow is $x\in\R^d$. For 5G-NIDD, each holdout leaves $K=8$ known classes and produces $d=88$ transformed features. For the CICIoT2023 category stress test, $K=7$ and $d=38$.

A source monitor outputs a probability vector $p(x)\in\Delta^{K-1}$, an attribution-proxy vector $\phi(x)\in\R^d$, and a normalized anomaly value $a(x)\in[0,1]$. The encoder $e$ produces a serialized message $m=e(x)$ with size $B(m)$ bytes. The known head returns
\begin{equation}
q_\theta(m)\in\Delta^{K-1},\qquad
\widehat{y}(m)=\argmax_{c\in\mathcal{Y}}q_{\theta,c}(m).
\label{eq:knownhead}
\end{equation}
The open-set head returns $S_\psi(m)\in[0,1]$ and rejects the observation when its calibrated p-value does not exceed $\alpha$.

\begin{table}[H]
\caption{Principal symbols.}
\label{tab:notation}
\centering
\begin{tabular}{ll}
\toprule
Symbol & Meaning \\
\midrule
$d$ & transformed predictor dimension ($88$ on 5G-NIDD; $38$ on CICIoT2023) \\
$K$ & number of known classes in one holdout ($8$ and $7$, respectively) \\
$m=e(x)$ & serialized evidence message \\
$B(m), B_{\max}$ & actual message bytes and allowed byte budget \\
$q_\theta(m)$ & known-class probability vector \\
$u,a,r$ & normalized uncertainty, anomaly, and prototype-residual evidence \\
$S_\psi(m)$ & scalar open-set score \\
$\tau$ & validation-calibrated rejection threshold \\
$\alpha$ & target known-traffic false-positive probability \\
\bottomrule
\end{tabular}
\end{table}

\subsection{Constrained multi-objective design}

Let $F_{\rm known}$ be known-class macro-F1, $R_{\rm unk}(\alpha)$ be held-out-family recall at false-positive budget $\alpha$, and $C$ be computational cost. The design problem is
\begin{equation}
\begin{aligned}
\max_{\theta,\psi,e}\quad&
\left(F_{\rm known}(\theta,e),\,R_{\rm unk}(\psi,e;\alpha)\right)\\
\text{subject to}\quad&
\FPR_{\rm known}(\psi,e)\leq\alpha,\quad
B(e)\leq B_{\max},\quad C(\theta,\psi,e)\leq C_{\max}.
\end{aligned}
\label{eq:multiobjective}
\end{equation}
Equation~\eqref{eq:multiobjective} is a constrained multi-objective formulation, not a scalar maximization of a vector. For a fixed application preference, a scalarized objective may be written
\begin{equation}
J_{\eta,\rho}=1-F_{\rm known}+\eta\left(1-R_{\rm unk}(\alpha)\right)
+\rho\,\frac{B(e)}{B_{\max}},
\label{eq:scalarization}
\end{equation}
subject to the same FPR and compute constraints. The experiments do not claim global Pareto optimality. They evaluate a finite encoding family and identify nondominated empirical operating points.
